\documentclass{article}

\usepackage{amsmath, amssymb, amsthm}
\usepackage[a4paper, margin=1in]{geometry}

\title{Selfridge Conjecture}
\author{sid}
\date{\today}

\newtheorem{theorem}{Theorem}
\newtheorem{lemma}{Lemma}
\newtheorem{example}{Example}

\newcommand{\Z}{\mathbb{Z}}
\newcommand{\bld}[1]{\textbf{#1}}
\newcommand{\itl}[1]{\textit{#1}}
\newcommand{\N}{\mathbb{N}}
\newcommand{\Q}{\mathbb{Q}}
\newcommand{\R}{\mathbb{R}}
\newcommand{\eps}{\varepsilon}
\newcommand{\poly}{\mathrm{poly}}
\newcommand{\polylog}{\mathrm{polylog}}
\newcommand{\bigO}{\mathcal{O}}

\begin{document}

\maketitle

John Selfridge has conjectured that if $p$ is an odd number, and  $p\equiv\pm2\bmod{5}$, then $p$ will be prime if both of the following hold:
\begin{enumerate}
  \item $2^{p-1}\equiv1\bmod{p}$
  \item $f_{p+1}\equiv0\bmod{p}$
\end{enumerate}
where $f_k$ is the $k$-th Fibonacci number.

The first condition is Fermat primality test using base 2.

\hrulefill

In general, if $p\equiv{a\bmod{(x^2+4)}}$ where $y^2\not\equiv{a\bmod{(x^2+4)}}\ \forall{y\in\N}$ ($a$ is a quadratic non-residue), then $p$ is a prime iff
\begin{enumerate}
  \item $2^{p-1}\equiv{1\bmod{p}}$
  \item $F_{p+1}(x)\equiv{0\bmod{p}}$
\end{enumerate}
$F_{k}(x)$ is the $k$-th Fibonacci polynomial at $x$.

$$
F_n(x) =
\begin{cases}
0, & \text{if } n = 0, \\
1, & \text{if } n = 1, \\
x F_{n-1}(x) + F_{n-2}(x), & \text{if } n \geq 2.
\end{cases}
$$

Selfridge, Pomerance and Wagstaff together offered \$620 for a counterexample or proof that there is none, with the prize now due from the Number Theory Foundation.

This eventually lead to Baillie-PSW primality test.

\end{document}